\clearpage
\begin{trapbox}
	\textbf{易错一｜横轴、纵轴不要认反}
\begin{center}\begin{tikzpicture}[font=\small,>=stealth,line join=round]
\draw[gray!40,rounded corners=3pt] (0,0) rectangle (14.2,8.3);
\node[text width=12.6cm,align=center] at (7.1,7.65) {关于 $x$ 轴：同一横坐标，上下高度相反};
\begin{scope}[shift={(4.3,4.1)},x=1.65cm,y=0.47000000000000003cm]
\draw[->,gray] (-2,0)--(2,0) node[right]{$x$};\draw[->,gray] (0,-5)--(0,5) node[above]{$y$};

\begin{scope}\clip (-2,-5) rectangle (2,5);
\draw[blue,thick,domain=-2:2,samples=180] plot(\x,{3*\x+2});
\draw[orange!85!black,thick,domain=-2:2,samples=180] plot(\x,{-3*\x-2});
\end{scope}
\fill[blue] (0,2) circle(2.2pt);\node[above left=6pt,blue] at (0,2) {$P$};
\fill[orange!85!black] (0,-2) circle(2.2pt);\node[below left=6pt,orange!85!black] at (0,-2) {$Q$};
\end{scope}
\node[anchor=west,align=left,text width=4.8cm] at (8.65,4.5) {蓝：$f(x)=3x+2$\\[8pt]橙：$g(x)=-3x-2$\\[8pt]$P(0,2)\to Q(0,-2)$\\[8pt]$g(x)=-f(x)$};
\node[text width=12.6cm,align=center] at (7.1,.65) {两幅图关于横轴相互对称；原函数本身不必关于横轴对称。};
\end{tikzpicture}\end{center}

\begin{description}[style=nextline, leftmargin=2.8em, labelsep=0.5em, itemsep=0.35em]
		\item[\textbf{\textcolor{CTrap}{易错点一（对称轴混淆）}}] 求关于 $y$ 轴对称时错误写为 $y=-f(x)$。
		\item[\textbf{\textcolor{CTrap}{正确理解（区分轴对称）}}] 关于 $x$ 轴对称：$y=-f(x)$；关于 $y$ 轴对称：$y=f(-x)$；两者符号作用位置不同。
		\item[\textbf{\textcolor{CTrap}{错因分析（变量混淆）}}] 未分清对称变换是对自变量还是因变量取反，导致公式用错。
	\end{description}

	\medskip
	\end{trapbox}
\clearpage
\begin{trapbox}
\textbf{易错二｜先检查反函数是否存在}
\begin{center}\begin{tikzpicture}[font=\small,>=stealth,line join=round]
\draw[gray!40,rounded corners=3pt] (0,0) rectangle (14.2,8.3);
\node[text width=12.6cm,align=center] at (7.1,7.65) {原函数不是一一对应：交换坐标后未必还是函数};
\begin{scope}[shift={(2.5666666666666664,3.3166666666666664)},x=0.7833333333333333cm,y=0.7833333333333333cm]
\draw[->,gray] (-2,0)--(4,0) node[right]{$x$};\draw[->,gray] (0,-2)--(0,4) node[above]{$y$};
\draw[violet,dashed] (-1.5,-1.5)--(3.6,3.6) node[above]{$y=x$};\draw[orange!85!black,thick,domain=-1.8:1.8,samples=150] plot({\x*\x},\x);\draw[gray,dashed] (1,-1.2)--(1,1.7);
\begin{scope}\clip (-2,-2) rectangle (4,4);
\draw[blue,thick,domain=-1.8:1.8,samples=180] plot(\x,{\x*\x});
\end{scope}
\fill[blue] (-1,1) circle(2.2pt);\node[above left=5pt,blue] at (-1,1) {$P$};
\fill[orange!85!black] (1,-1) circle(2.2pt);\node[below right=5pt,orange!85!black] at (1,-1) {$Q$};
\fill[orange!85!black] (1,1) circle(2.2pt);\node[above right=5pt,orange!85!black] at (1,1) {$(1,1)$};
\end{scope}
\node[anchor=west,align=left,text width=4.8cm] at (8.65,4.5) {蓝：$y=x^2$\\对称后是 $x=y^2$\\[8pt]当 $x=1$：\\同时有 $y=1,-1$\\[8pt]同一输入有两个输出，\\不能直接写成反函数。};
\node[text width=12.6cm,align=center] at (7.1,.65) {两轴同单位；原函数限于 $x\ge0$ 后，反函数为 $y=\sqrt{x}$。};
\end{tikzpicture}\end{center}

\begin{description}[style=nextline, leftmargin=2.8em, labelsep=0.5em, itemsep=0.35em]
		\item[\textbf{\textcolor{CTrap}{易错点二（反函数未显化）}}] 写出 $x=f(y)$ 后，不检查同一个 $x$ 是否会对应多个 $y$，就断言存在反函数。
		\item[\textbf{\textcolor{CTrap}{正确理解（显化反函数）}}] 交换坐标所得关系 $x=f(y)$ 本身是正确的；能否表示为反函数，关键在于 $f$ 是否一一对应。题目要求显式解析式时，再解出 $y$。
		\item[\textbf{\textcolor{CTrap}{错因分析（函数与方程）}}] 函数要求一个输入恰好对应一个输出；不是单凭纵坐标有没有单独写在等号左边来判断。
	\end{description}

	\medskip
	\end{trapbox}
\clearpage
\begin{trapbox}
\textbf{易错三｜完整定义域也要反射}
\begin{center}\begin{tikzpicture}[font=\small,>=stealth,line join=round]
\draw[gray!40,rounded corners=3pt] (0,0) rectangle (14.2,8.3);
\node[text width=12.6cm,align=center] at (7.1,7.65) {对称后定义域也要反射：不是一律沿用原范围};
\begin{scope}[shift={(2.0999999999999996,2.4214285714285717)},x=1.0999999999999999cm,y=1.342857142857143cm]
\draw[->,gray] (-1,0)--(5,0) node[right]{$x$};\draw[->,gray] (0,-0.5)--(0,3) node[above]{$y$};
\draw[violet,dashed] (2,-.3)--(2,2.7) node[above]{$x=2$};
\begin{scope}\clip (-1,-0.5) rectangle (5,3);
\draw[blue,thick,domain=0:5,samples=180] plot(\x,{sqrt(\x)});
\draw[orange!85!black,thick,domain=-1:4,samples=180] plot(\x,{sqrt(4-\x)});
\end{scope}
\fill[blue] (1,1) circle(2.2pt);\node[above left=5pt,blue] at (1,1) {$P$};
\fill[orange!85!black] (3,1) circle(2.2pt);\node[above right=5pt,orange!85!black] at (3,1) {$Q$};
\fill[blue] (0,0) circle(2.2pt);\node[below=5pt,blue] at (0,0) {$0$};
\fill[orange!85!black] (4,0) circle(2.2pt);\node[below=5pt,orange!85!black] at (4,0) {$4$};
\end{scope}
\node[anchor=west,align=left,text width=4.8cm] at (8.65,4.5) {蓝：$f(x)=\sqrt{x}$\\定义域 $[0,+\infty)$\\[8pt]关于 $x=2$ 反射\\[8pt]橙：$g(x)=\sqrt{4-x}$\\定义域 $(-\infty,4]$\\[8pt]$1\leftrightarrow3$，高度均为 $1$。};
\node[text width=12.6cm,align=center] at (7.1,.65) {截取窗口内曲线；蓝图向右延伸，橙图向左延伸，整幅定义域必须一起对应。};
\end{tikzpicture}\end{center}

\begin{description}[style=nextline, leftmargin=2.8em, labelsep=0.5em, itemsep=0.35em]
		\item[\textbf{\textcolor{CTrap}{易错点三（定义域忽视）}}] 对称变换后直接使用原函数定义域，未考虑新函数定义域变化。
		\item[\textbf{\textcolor{CTrap}{正确理解（重新确定定义域）}}] 新函数的定义域应为对称后图像上点的横坐标取值范围，需根据原函数值域或对称情况重新确定。例如 $f(x)=\sqrt{x}$ 关于 $y$ 轴对称后定义域变为 $x\le 0$。
		\item[\textbf{\textcolor{CTrap}{错因分析（条件忽视）}}] 只关注解析式变换，忽略了函数定义域是函数的重要组成部分，对称后横坐标范围改变。
	\end{description}
\end{trapbox}
\clearpage
\begin{trapbox}
\textbf{易错四｜把两图相互对称的轴写成 $(a+b)/2$}
\begin{center}\begin{tikzpicture}[font=\small,>=stealth,line join=round]
\draw[gray!40,rounded corners=3pt] (0,0) rectangle (14.2,8.3);
\node[text width=12.6cm,align=center] at (7.1,7.65) {把轴误写成 $x=3$：最低点的对应点落不到另一图上};
\begin{scope}[shift={(2.32,2.3375)},x=0.6599999999999999cm,y=0.5875cm]
\draw[->,gray] (-2,0)--(8,0) node[right]{$x$};\draw[->,gray] (0,-1)--(0,7) node[above]{$y$};
\draw[red!70!black,dashed] (3,-.7)--(3,6.5);\node[red!70!black,font=\footnotesize] at (3.4,7.8){$x=3$};\draw[red!70!black,dashed] (-1,0)--(7,0);\draw[red!70!black] (7,0) circle(2.7pt);\node[below=7pt,red!70!black] at (7,0){$(7,0)$};
\begin{scope}\clip (-2,-1) rectangle (8,7);
\draw[blue,thick,domain=-2:8,samples=180] plot(\x,{(1+\x)^2});
\draw[orange!85!black,thick,domain=-2:8,samples=180] plot(\x,{(5-\x)^2});
\end{scope}
\fill[blue] (-1,0) circle(2.2pt);\node[above=10pt,blue] at (-1,0) {$P$};
\fill[orange!85!black] (7,4) circle(2.2pt);\node[above left=5pt,orange!85!black] at (7,4) {$(7,4)$};
\end{scope}
\node[anchor=west,align=left,text width=4.8cm] at (8.65,4.5) {仍取 $F(x)=(x+1)^2$\\$G(x)=(5-x)^2$\\[8pt]若错用轴 $x=3$：\\$(-1,0)\to(7,0)$\\[8pt]但 $G(7)=4\ne0$。\\正确轴仍为 $x=2$。};
\node[text width=12.6cm,align=center] at (7.1,.65) {红色空心点是错误候选的镜像点，不是橙色函数图像上的点。};
\end{tikzpicture}\end{center}


\textbf{错误：}$y=f(a+x)$ 与 $y=f(b-x)$ 的轴是 $x=(a+b)/2$。
\textbf{反例：}图中 $a=1,b=5$，错误候选 $x=3$ 会把蓝图最低点 $(-1,0)$ 送到 $(7,0)$，但橙图在 $x=7$ 的高度为 $4$，不是 $0$。
\textbf{正确替代：}用 $F(2c-x)=G(x)$ 检查，得到 $a+2c=b$，所以 $c=(b-a)/2$。

\end{trapbox}
\clearpage
\begin{trapbox}
\textbf{易错五｜混淆相互对称与自身对称}
\begin{center}\begin{tikzpicture}[font=\small,>=stealth,line join=round]
\draw[gray!40,rounded corners=3pt] (0,0) rectangle (14.2,8.3);
\node[text width=12.6cm,align=center] at (7.1,7.65) {关于 $x$ 轴：同一横坐标，上下高度相反};
\begin{scope}[shift={(4.3,4.1)},x=1.65cm,y=0.47000000000000003cm]
\draw[->,gray] (-2,0)--(2,0) node[right]{$x$};\draw[->,gray] (0,-5)--(0,5) node[above]{$y$};

\begin{scope}\clip (-2,-5) rectangle (2,5);
\draw[blue,thick,domain=-2:2,samples=180] plot(\x,{3*\x+2});
\draw[orange!85!black,thick,domain=-2:2,samples=180] plot(\x,{-3*\x-2});
\end{scope}
\fill[blue] (0,2) circle(2.2pt);\node[above left=6pt,blue] at (0,2) {$P$};
\fill[orange!85!black] (0,-2) circle(2.2pt);\node[below left=6pt,orange!85!black] at (0,-2) {$Q$};
\end{scope}
\node[anchor=west,align=left,text width=4.8cm] at (8.65,4.5) {蓝：$f(x)=3x+2$\\[8pt]橙：$g(x)=-3x-2$\\[8pt]$P(0,2)\to Q(0,-2)$\\[8pt]$g(x)=-f(x)$};
\node[text width=12.6cm,align=center] at (7.1,.65) {两幅图关于横轴相互对称；原函数本身不必关于横轴对称。};
\end{tikzpicture}\end{center}


\textbf{错误：}两图 $y=f(x)$ 与 $y=-f(x)$ 关于 $x$ 轴相互对称，就说 $f$ 的图像本身关于 $x$ 轴对称。
\textbf{为什么错：}图中蓝线 $3x+2$ 与橙线 $-3x-2$ 是两条不同的线。蓝点 $(0,2)$ 的镜像 $(0,-2)$ 在橙线上，不在蓝线上。
\textbf{正确替代：}先区分是哪两幅图。仅当反射后的整个图像与原图重合时，才是自身对称；关于竖直轴或点的自身判定分别见 F020、F021。

\end{trapbox}
